\section{Machine Learning Model for Reduction Propagation}
\label{sec:reduction-model}
The seed model selects identities from which a reduction can proceed.
Once these equations are available, the reduction model selects the next
row operation from a supplied list of allowed steps. We use the current
coefficient matrix and its integral labels to describe each choice, and
the model returns a probability for every listed step. The step with the
largest probability gives the proposed next operation. Exact arithmetic
then updates the matrix, which supplies the input for the next choice.
This section describes the model and its training on the reductions
constructed in Section~\ref{sec:integrals-generation}.

\subsection{Current equations and allowed reduction steps}
Let $A^{(n)}$ be the coefficient matrix after $n$ reduction steps. A column
$j$ corresponds to an integral $I_{\boldsymbol a^{(j)}}$, and a row contains
its coefficients in an identity or in the target expression currently
being reduced. The chosen master basis specifies the columns protected
from elimination. An action $a=(q,r,k)$ uses source row $q$ to eliminate
column $k$ from destination row $r$:
\begin{equation}
 A^{(n+1)}_{rj}
 =A^{(n)}_{rj}
  -\frac{A^{(n)}_{rk}}{A^{(n)}_{qk}}A^{(n)}_{qj},
 \qquad q\ne r,\quad A^{(n)}_{qk}\ne0,\quad A^{(n)}_{rk}\ne0.
 \label{eq:reduction-action}
\end{equation}
All other rows remain unchanged during this elementary step. If both
rows are identities, this is an elementary row operation. If row $r$
represents a target, row $q$ must be a zero identity so that the target
retains its value. The solver supplies the allowed action set
$\mathcal A_n$, enforcing its integral ordering, master basis and row
roles. The model ranks this finite list; completeness of the list remains
the responsibility of the solver.

The stored examples describe a particular instance of this matrix.
They contain one changing target row and the source identities available
for eliminating its largest remaining nonmaster integral under
Eq.~\eqref{eq:target-reference-order}. Each candidate specifies one of
these identities, the target row and that integral. After a choice, the
source identities remain fixed and the target row changes. The prediction
interface also accepts several changing rows, while the present training
data cover the stored target-row procedure. Section~\ref{sec:live-matrix} tests the frozen policy on Kira's evolving
selected IBP matrix using admissible actions that retain its integral
ordering~\cite{Kira3}. This test changes the inference states while retaining
the target-row training described here.

\subsection{Information supplied to the model}
For each allowed action, we construct a vector $\boldsymbol f_a$ from
the current matrix, the pivot integral and the family definition.
The pivot description includes its indices, sector membership, and
the values of $r,s,t,n_{\rm dot}$ in Eq.~\eqref{eq:complexities}.
Descriptions of the source and destination rows give their nonzero
counts, master content and the mean and largest integral complexities.
A matrix summary gives its number of rows, columns and nonzero entries.
The public family definition supplies the numbers of loops and
kinematic parameters together with the propagators and scalar-product
relations. We encode their written form through the same fixed signed
token counts described in Section~\ref{sec:seed-model-architecture},
using 32 components here. This encoding admits collisions between
different groups of symbols. We transform nonnegative counts with
$\log(1+v)$, where $v$ is the count, and each signed index $a_i$ with
$\operatorname{sgn}(a_i)\log(1+|a_i|)$. The representation allows up to
16 indices and records which positions belong to the family.
These definitions give 132 features for each candidate.

The current records supply coefficients in $\mathbb F_P$, the field
of integers modulo a prime $P$. We use these coefficients to calculate
the overlap of each source row with its destination, the number of
terms that the proposed operation would cancel or introduce, and its
immediate arithmetic count. For a source with $n_q$ nonzero entries,
this count is $1+2n_q$: one division and $n_q$ multiplications and
subtractions, using the convention in
Section~\ref{sec:integrals-generation}. These quantities require one
sparse trial row update per action. They use only the current matrix
and do not inspect later reduction steps or reference master
coefficients.
A symbolic coefficient can vanish at the chosen sample, so the model
uses counts recomputed from the sampled rows. The data audit retains
the distinction from the original symbolic support counts. Structural
counts and integral descriptions carry this information into the model.
The numerical size of a modular residue does not represent the
magnitude of a real coefficient.

The representation compresses the supplied matrix into these features,
so different states can have the same description. It includes local
cancellation information and a summary of the candidate set. These
quantities describe the current choices without enumerating complete
reduction paths. Forming these features and evaluating the model
contribute to the eventual runtime
and must be included when measuring a reduction speedup.

\subsection{Action probabilities and confidence}
A shared map $\phi_\theta$ converts each candidate vector into a
32-component vector,
\begin{equation}
 \boldsymbol h_a=\tanh(W_1\boldsymbol f_a+\boldsymbol b_1),\qquad
 \boldsymbol u=\frac{1}{|\mathcal A_n|}
                    \sum_{a\in\mathcal A_n}\boldsymbol h_a,\qquad
 \boldsymbol v=\max_{a\in\mathcal A_n}\boldsymbol h_a .
 \label{eq:reduction-encoding}
\end{equation}
Here $\theta$ denotes the adjustable parameters, $W_1$ is a learned
matrix, $\boldsymbol b_1$ is a learned offset, and the maximum acts
on each component separately. Sharing the map and combining candidates
through averages and maxima makes the scores independent of their
position in the supplied list~\cite{DeepSets}. The model joins
$\boldsymbol h_a,\boldsymbol u,\boldsymbol v$ and applies a second
map with 64 hidden components to obtain a scalar score $g_\theta(a)$.
For a fixed family, master basis and finite-field sample, it returns
\begin{equation}
 \begin{aligned}
 p_\theta(a\mid A^{(n)},\mathcal A_n)
 &=\frac{\exp g_\theta(a)}
        {\sum_{b\in\mathcal A_n}\exp g_\theta(b)},\\
 \widehat a_n&=\mathop{\rm arg\,max}_{a\in\mathcal A_n}p_\theta(a),
 \qquad q_n=p_\theta(\widehat a_n).
 \end{aligned}
 \label{eq:reduction-choice}
\end{equation}
Thus the probabilities sum to one over the supplied alternatives,
and $q_n$ is the reported confidence for the selected step.
With one allowed step the probability is one by construction, which
provides no evidence of learned discrimination. An empty candidate
list returns no proposed operation; the solver must determine whether
the reduction is complete or further identities are required.

Training gives these probabilities an interpretation as agreement
with the recorded next choice, conditional on the supplied state and
actions. Several choices can give valid reductions, and an alternative
choice may even lead to a shorter sequence. The reference paths establish
valid reductions and their recorded costs; they do not establish global
minima. Section~\ref{sec:reduction-test} compares this confidence with agreement
on reserved choices. Confidence does not certify exactness, minimal
cost or the success of a complete reduction. Equation~\eqref{eq:reduction-action}
and the solver's checks retain responsibility for the algebra.

\subsection{Training examples and separation from testing}
We retain the same grouped family division used for the seed model:
11,963 families with 1,610,890 targets supply training records, while
5,137 families with 690,382 targets remain reserved for testing.
Every state from one target stays with that target's family.
We read the row before each recorded step and its complete stored list
of admissible source equations. The chosen equation supplies the
reference action $a_n^\star$. The row after the step and the stored
operation counts check the data preparation but remain outside
the predictive inputs. The target's construction history, final master
coefficients and future choices also remain outside those inputs.

For the set $\mathcal D_{\rm choice}$ of training states with at least
two allowed actions, the objective is
\begin{equation}
 \mathcal L_{\rm red}(\theta)
 =-\frac{1}{|\mathcal D_{\rm choice}|}
   \sum_{(A,\mathcal A,a^\star)\in\mathcal D_{\rm choice}}
           \log p_\theta(a^\star\mid A,\mathcal A).
 \label{eq:reduction-loss}
\end{equation}
Each such state has equal weight. States with one admissible action
have zero loss and remain in the data audit. The recorded reference
choice is the training label; absence of a label on another action does
not establish that its reduction is invalid.

We use AdamW~\cite{AdamW} with a learning rate of $10^{-3}$, weight decay
$10^{-4}$, batches of 256 states and a gradient norm limit of one.
We continued the saved expanded-data model and optimizer for twenty
complete passes over the enlarged training states. The new phase
retains the architecture and optimizer settings while adding 1,250
mass profiles at every represented loop order. Preparation replays every selected row
operation with exact modular arithmetic and verifies its final
recorded master row. Training shuffles states within bounded
groups of families and retains complete resume states.
The reserved families remain outside parameter updates and
confidence fitting. Earlier project stages evaluated the preceding
held-out families; the current experiment also reserves the newly
assigned mass profiles at every added loop order.

\subsection{Completed training and present scope}
The expanded training set contains 14,684,059 recorded states from
1,610,890 targets in 11,963 families. Of these states,
5,840,122 supply at least two allowed actions and
8,843,937 supply one action. The records contain
26,775,999 candidate appearances. The unchanged model
uses 132 features per candidate and has 10,529 adjustable parameters.

Preparation replays every reference row operation, checks continuity
between steps and verifies the final master support. Symbolic
coefficients can vanish at the sampled kinematic point, so the model
uses the support of the actual sampled matrix. The audit records
these differences without changing the exact symbolic references.

We continued the preceding expanded-data model and its AdamW
optimizer for twenty complete passes over the enlarged data. Each new pass
uses every training state with an action choice; a forced step has
probability one and supplies no ranking decision. The new phase
contains 116,802,440 choice-state presentations and
463,763 parameter updates. The preserved lineage contains thirty
passes over the original smaller collection and twenty passes over
the preceding 2,100-family collection. The latest twenty passes
use the enlarged collection across all represented loop orders.

Mean training cross-entropy changed from 0.3823373 to
0.3804633 during this phase, while agreement with the recorded
action changed from 99.841\% to 99.895\%.
These values describe losses and choices encountered during
parameter updates. Section~\ref{sec:reduction-test} measures the
frozen parameters on reserved families.

The implementation retains the original architecture, learning rate,
gradient bound and state-batch limits. Training uses four CPU threads,
while inference uses one. Parallel preparation reproduces the saved
training features and labels exactly. Checks verify equal weights
in the inference and resumable checkpoints, complete state coverage
and the preserved optimizer lineage. The published implementation
retains the inputs, learned parameters and exact reference
records~\cite{ReductionModelData}.

The training domain still consists of one changing target row and a
fixed source-equation pool. The new records extend this domain
through twelve loops for the coupled vacuum families specified
above. They provide evidence about these allowed reduction steps;
general sparse elimination inside Kira requires a separate
integration and benchmark.

